\begin{textsample}{POS Dim 3 – ba – Score 20.00 – ba/ba\_inf\_10.txt}  \label{ex:f3_pos_201}
7.10.
Campos de Experiência no Currículo da Educação \textbf{Infantil} Como pensar uma Base Nacional Comum Curricular para todos os níveis de educação e ensino sem perder a especificidade da Educação \textbf{Infantil} e, aquilo que Carlina Rinaldi chamou de “ Currículo Emergente ” ( 1999 ), com a proposta de diferenciar -se de um planejamento como um método de trabalho que estabelece de antemão objetivos educacionais gerais e específicos?
Cabe então questionar o formato do currículo escolar organizado por disciplinas que, muitas vezes, é trazido como modelo para Educação \textbf{Infantil}, de modo que possamos pensar novas formas de lidar com os saberes, materiais, tempos e espaços educacionais específicos da Educação \textbf{Infantil} para as \textbf{crianças} \textbf{pequenas}.
Por tal razão, destacamos a necessidade de refletir sobre os campos de experiência no contexto da educação da infância e suas contribuições para pensar o processo de construção de conhecimentos, para construir um processo educativo que considere as trocas entre as \textbf{crianças} e entre \textbf{adultos} e \textbf{crianças}.
Em síntese, buscar contribuir para um processo educativo que tem na \textbf{criança} a sua centralidade.
Nesse caso, a dúvida e a fascinação são fatores muito bem-vindos, juntamente com a investigação, a descoberta e a invenção.
Com base no art. 9º das Diretrizes Curriculares Nacionais para Educação \textbf{Infantil}, abre -se a proposição dos Campos de Experiências como proposta norteadora para a organização curricular da Educação \textbf{Infantil}.
Essa abordagem considera a experiência da \textbf{criança} como sujeito que age, cria e produz cultura, visão que supera a lógica da \textbf{criança} como mera receptora dos conhecimentos.
Os campos de experiências colocam o fazer e o agir da \textbf{criança} no centro do saber, além de articular os direitos de aprendizagens da \textbf{criança}, que destacam seis princípios básicos, a saber : conviver, \textbf{brincar}, participar, explorar, expressar e conhecer -se. \textbf{ESCUTA}, FALA, PENSAMENTO E IMAGINAÇÃO Como um dispositivo estruturante e propositivo de Currículos, os campos de experiência trazem transformações significativas para a qualificação dos processos formacionais, na medida em que evitam o abstracionismo e a insularidade da tradição disciplinar.
No que se refere à Educação \textbf{Infantil}, dá centralidade à aprendizagem da \textbf{criança} e seu processo formacional, compreendidos como singularidades experienciais.
Nesses termos, aprendizagem e formação como experiências valoradas estão apartadas de qualquer perspectiva exterodeterminante da aprendizagem e da formação.
Dessa forma, processos de autonomização e de autorização por parte da \textbf{criança} implicam em aprender e desenvolver saberes que envolvem exprimir sentimentos e elaborar estratégias, expressar pontos de vista e aventurar -se honestamente em definições de situação e, com isso, na construção da afirmação atitudinal.
É com esse conjunto de dispositivos e capacidades adquiridas que as \textbf{crianças}, ao \textbf{brincar}, ao exercitar sua \textbf{curiosidade}, ao questionar, ao explorar, ao narrar, ao escutar, produzem descritibilidades, inteligibilidades, analisibilidades e sistematicidades, compreendidas como competências que as ajudam a questionar e criar hipóteses sobre o mundo e seus sistemas simbólicos, de uma forma sempre ativa, mesmo que nessas itinerâncias produzam erros.
A construção de processos de autonomização, de autorização e autorias \textbf{infantis} se constitui em aventuras pensadas, eivadas de intensas atividades cognitivas, desiderativas, sociais e simbólicas. 21 BRASIL.
Conselho Nacional de Educação ; Câmera de Educação Básica.
Parecer nº 7, de 7 de abril de 2010.
Diretrizes Curriculares Nacionais Gerais para a Educação Básica.
Diário Oficial da União, Brasília, 9 de julho de 2010, Seção 1, p. 10.
Disponível em :.
Acesso em : 19 set. 2018.
A aprendizagem é, acima de tudo, ação simbólica, deslocamento e alteração do ser \textbf{criança} na sua totalidade.
Enquanto processo, ao mesmo tempo individual, relacional e simbólico, a aprendizagem, enquanto saber, saber-fazer e saber-ser, é fundamentalmente uma atividade experiencial.
Assim, mediada por procedimentos pedagógicos na escola, bem como em outros tempos e espaços, se dá na experiência do ser que aprende e sua singularidade.
Nesses termos, nos processos formacionais das \textbf{crianças}, atitudes de pesquisa colaboram para que \textbf{crianças} acordem as fontes dos saberes, experimentem intensamente, pensem e expressem seus pensamentos, dialoguem tentando esclarecê -los, tendo a narrativa um lugar de destaque nesse processo pedagógico.
Justifica -se aqui não só valorizar e valorar a experiência e os saberes, mas os \textbf{acolher} na sua plenitude, dialogando com sua pertinência e relevância, assim como problematizando -os quando necessário.
É preciso que o educador esteja atento para o fato de que nem toda aprendizagem é boa.
Importa criar uma ambiência onde o diálogo honesto, acolhedor, crítico possa valorizar e valorar as aprendizagens, tornando -as potencialmente formadoras para as experiências \textbf{infantis} e para o próprio educador, até porque \textbf{crianças} os fazem aprender também.
O encontro mediado entre educadores e \textbf{crianças} é mutuamente formacional.
Com esses argumentos, faz -se necessário afirmar que trabalhar com os campos de experiência na Educação \textbf{Infantil} não significa cair no populismo pedagógico de que a mediação do educador deixa de ser propositiva ou mesmo crítica.
Acolhedora, solidária, pedagogicamente preparada, essa mediação é, acima de tudo, dialógica, trata -se de um encontro experiencial transversalizado, por intenções pedagógicas compromissadas com aprendizagens formacionais, portanto, qualificadas.
Implicar nos processos pedagógicos as explorações \textbf{infantis}, suas \textbf{curiosidades}, proposições e descobertas, não significa o educador se ausentar propositivamente do processo de qualificação pleiteada pelo Currículo.
Sabe -se, a fortiori, que, em realidade, quem constrói os saberes e suas formas de ação são as \textbf{crianças}, mas num mundo intensamente relacional.
Jogos e \textbf{brincadeiras}, como centralidades \textbf{lúdicas}, constroem e oportunizam descobertas e aprendizagens.
Nesses termos, é preciso cuidar das ambiências pedagógicas e seus dispositivos, para que essa condição e direito da infância seja vivido plenamente e potencialize aprendizagens formacionais.
Cada campo de experiência oferece um conjunto de objetos, situações, imagens e linguagens, relacionados aos sistemas simbólicos da nossa cultura, capazes de evocar, estimular, acompanhar \textbf{progressivamente} aprendizagens.
Pode -se falar que os campos de experiência como organizadores curriculares necessitam de uma ecologia educacional para que seu potencial formacional possa emergir.
Essa é uma recomendação fulcral para a gestão curricular da Educação \textbf{Infantil}, que implica conceber, organizar e disponibilizar ambiências saudáveis e artefatos culturais pedagogicamente relevantes.
Campos de Experiência como dispositivo organizador de um currículo significam, acima de tudo, colocar a experiência da \textbf{criança} no centro das atenções a partir de disponibilidades interativas que se disponibilizem a trabalhar com essas experiência, e não sobre elas ou contra elas.
É compreendê -las sempre como fonte de saberes de possibilidades formacionais.
É por meio da experiência que se dá sentido ao mundo, às nossas atividades e realizações.
É com as experiências cotidianas que as \textbf{crianças} buscam e constroem suas políticas de sentido, ou seja : fazem opções, compreendem o mundo e tomam decisões.
Por outro lado, faz -se necessário, no encontro entre a experiência docente e da \textbf{criança} em formação, intercriticizar seus sentidos para que se possa valorizar e valorar seus conteúdos formacionais.
Cabe nesse contexto reinventar a docência para trabalhar com a complexidade e os desafios das diversas experiências num âmbito pedagógico estruturado por ações-reflexões-ações implicadas à qualificação dos processos formacionais.
O eu, o outro, o nós Na infância, a construção dos processos identitários, realizada por meio de inúmeras formas de observação e de indagações que envolvem as pessoas, eventos, tradições familiares, culturas em que as \textbf{crianças} estão inseridas desde os seus primeiros dias de vida, efetiva comparações, assim como processos de inclusão e exclusão.
Neste processo, a \textbf{criança} \textbf{interessa} -se pela sua existência e das outras pessoas, entre outros seres.
Presentes nessas questões surgem, mais tarde, perguntas sobre Deus, a vida e a morte, a tristeza e a felicidade.
Como muitas vezes se tratam de questões complexas, faz -se necessário uma \textbf{escuta} compreensiva e refinada para que, de acordo com sua capacidade cognitiva e afetiva, respostas formacionais cuidadosas sejam construídas pelos seus educadores. \textbf{Crianças} constroem seus processos identitários convivendo e dialogando consigo e com os outros.
Experimentam estados de humor e, com isso, aprendem a expressá -los, em busca de apoio, \textbf{cuidado}, proteção e interação qualificada.
É aqui que muitos dos seus direitos deverão ser exercidos, tendo como guardiões o Estado e a família, entre outras instituições educacionais, meios pelos quais também aprendem sobre seus deveres.
Da perspectiva da infância, esta é a condição de uma passagem evolutiva, importante, na superação gradual do seu egocentrismo, entrando em cena, com importantes aprendizagens sociais.
É aí que a escola tem um papel social fulcral no processo de ampliação dos processos de socialização e, portanto, de ampliação, também, do processo identitário da \textbf{criança}.
É aí, também, que acontece uma diferenciação significativa da qualificação da convivência, a partir de valores vinculados à solidariedade, à reciprocidade e ao respeito dos direitos e deveres de si próprio e dos outros.
O outro começa a surgir como fonte de possibilidades e limites, assim como valores democráticos importantes a serem exercidos pela cidadania.
Corpo, \textbf{gestos} e movimentos Trata -se da importância da tomada de consciência do corpo e da compreensão de que toda a nossa vida passa pela condição corporal, suas amplas e, às vezes, insondáveis experiências.
Movimentos, \textbf{gestos}, como correr e pular, produzem bem-estar e \textbf{equilíbrio} psicofísico.
Sensações e emoções são produzidas e necessariamente passam pelo corpo.
Aliás, toda e qualquer experiência passa pelo corpo como lugar de estados de ser.
Relaxamento, tensão, controle de \textbf{gestos}, limites e possibilidades físicas implicam aprendizagens importantes para a luta pela qualificação da vida.
As \textbf{crianças} jogam com o próprio corpo, comunicam -se e exprimem -se com a \textbf{mímica}.
As experiências motoras permitem integrar as diferentes linguagens.
Jogos que impliquem a psicomotricidade fina e ampla constroem aprendizagens importantes, assim como satisfação e saúde.
Nesses termos, é de suma importância que uma arquitetura de \textbf{prédios} para Educação \textbf{Infantil} tenha consciência da importância do planejamento dos espaços para que a especificidade pedagógica da Educação \textbf{Infantil} tenha lugar.
É assim que a Educação \textbf{Infantil} e seus espaços adequados possibilitam a expressão e a comunicação pelo corpo, assim como as diversas expressões artísticas pelas quais a \textbf{criança} aprende a se movimentar em diversos e complexos tempos e espaços da vida.
Conhecer e cuidar do seu corpo, assim como compreender que o corpo do outro merece \textbf{cuidado} e respeito, é parte de uma formação valorosa e valorada do ser da \textbf{criança}.
Traços, \textbf{sons}, cores e formas Em geral, as \textbf{crianças} se encontram nas Artes com uma facilidade impressionante.
É também por isso que a Educação \textbf{Infantil} encontra nas Artes potenciais significativos para a formação da \textbf{criança}.
Além da criatividade, a Arte implica emoções, imaginação, sensibilidade e autoria artística.
Arte e diferença são entretecimentos que criam singularidades incessantes, ao mesmo tempo em que elaboram experiências formacionais, singularizantes, porque vivem da e na criação.
A experimentação de materiais e linguagens como a música, a dramatização, os \textbf{sons}, elaborações gráfico-pictóricas, bem como a criação e experimentação de mídias, implicam atitudes de pesquisa e um prazer singular nos processos de aprendizagem.
Em termos contemporâneos, as experiências com a transmídia inserem as \textbf{crianças} em verdadeiros cenários de experimentações que, constantemente, as colocam no devir artístico, numa cultura contemporânea, a cibercultura, por exemplo, que lhes desafiam prazerosamente, tanto individual quanto coletivamente.
É aí que a fruição e a invenção elevam a imaginação \textbf{infantil} a mundos a serem explorados e observados com significativos potenciais formativos.
Essas possibilidades estão no museu, no cinema, no circo, nas instalações, nos espetáculos de rua, no teatro, nos eventos musicais, na televisão, no digital etc.
Patrimônio artístico que se encontra como possíveis patrimônios formacionais da \textbf{criança}. \textbf{Escuta}, fala, pensamento e imaginação A linguagem não só expressa o pensamento, ela é generativa, ou seja, criativa, assim como mobiliza o desenvolvimento cognitivo.
Nesses termos, implica -se o comunicar e a complexidade do compreender.
No processo cultural, emerge a língua como fenômeno identitário da \textbf{criança}, bem como um patrimônio cultural afirmativo da singularidade de um povo.
No encontro entre línguas, amplia -se a compreensão cultural da \textbf{criança}.
As \textbf{crianças}, no encontro com a escola, já possuem um repertório linguístico rico.
Na escola, esse repertório se diferencia na medida em que processos de socialização e outros são vividos a partir de aprendizagens mais ampliadas e dirigidas para competências coletivas, socialmente referenciadas.
Ouvindo histórias e contos, confrontando pontos de vista, experimentando jogos e atividades mais formalizadas, interagindo com \textbf{adultos} e \textbf{colegas}, criando jogos com a língua, as \textbf{crianças} exploram possibilidades e produzem outras formas de comunicação prazerosa.
Num ambiente estimulante e acolhedor, o encontro com a língua escrita por meio de livros ilustrados, de mensagens, de orientações escritas, da convergência de mídias, a \textbf{criança} amplia o léxico e amplia, também, a aproximação com \textbf{sons}, palavras e frases corretas, e, com isso, vai experimentando o prazer da comunicação por meio da língua.
Assim, os sentidos enriquecem e outros mundos podem ser compreendidos nas suas aproximações e diferenças.
Espaços, tempos, quantidades, relações e transformações A \textbf{curiosidade} da \textbf{criança} é uma das atitudes que favorecem a sua inserção no querer saber no que se refere aos diversos fenômenos da Natureza.
Convidá -la a refletir sobre conceitos científicos e matemáticos pode ser pedagogicamente um prolongamento de sua condição para questionar quase tudo.
Nesses termos, ao observar os fenômenos que os cercam, tentam compreender experimentando e observando suas mudanças.
Neste mesmo fluxo de \textbf{curiosidades} sobre si, sobre outros seres, assim como sobre algumas lógicas da tradição matemática contidas em histórias e jogos matemáticos, ampliam a compreensão do mundo em que vivem.
Na vontade de experimentar e questionar o mundo em que vive e as informações a que têm acesso, a atitude de pesquisa num sentido amplo já está presente no modo de ser \textbf{criança}.
Nesse sentido, podemos considerar a \textbf{criança} um ser que gosta de exercitar a exploração e o mundo para compreendê -lo.
Como mediadores pedagógicos desse processo, o professor e seus auxiliares e \textbf{adultos} próximos ajudam as \textbf{crianças} a ampliar e complexificar suas compreensões, formulando questões explicitativas e oferecendo pistas, à medida que os \textbf{pequenos} exploram objetos, materiais e manifestações da Natureza.
Nesse processo, os professores também vão mediando formas pedagógicas mais estruturadas e sistematizadas de oportunizar na \textbf{criança} compreensões lógicas.
Pegando e sentindo o movimento dos objetos, as \textbf{crianças} sentem sua duração e movimento, aprendendo assim a organizá -los a partir da sua interferência na realidade.
Nesse caminho de experimentações, tomam gosto pela construção e reconstrução dos mesmos e, em algumas situações, na montagem e desmontagem das suas configurações.
No processo de observação e experimentação dessas ações, aprimoram seus \textbf{gestos} e suas compreensões.
Aprendem a fazer perguntas, a dar e a pedir explicações, a se deixar convencer pelos pontos de vista dos outros, a não se desencorajar se suas ideias não resultam apropriadas.
Podem, portanto, dar início a uma itinerância de conhecimentos mais estruturados, em que exploram as potencialidades da linguagem para se exprimir e usar símbolos para representar sentidos e significados.
Envoltas nessas experiências, as \textbf{crianças} vão criando questões e hipóteses de como funcionam e funcionariam os objetos com os quais trabalham e convivem cotidianamente.
É nesse momento que, pedagogicamente, passa a ser importante trabalhar a própria estruturação, desenvolvimento e funcionalidade do próprio corpo, bem como as relações possíveis com outros corpos de seres vivos e suas formas de vida.
Com esse movimento de \textbf{curiosidades}, experimentações e compreensões em processo de ampliação, inseri -las na possibilidade de entendimento de como funcionam os fenômenos invisíveis da Natureza e do seu corpo aguça a \textbf{curiosidade} e amplia os instrumentos cognitivos de compreensão das realidades que não estão ao alcance das suas observações diretas.
Com isso, oportunizam -se os exercícios cognitivos que vão se afastando de um pensar, colando ao que é observável.
Acessar instrumentos que permitem visibilizar o invísivel ajuda a \textbf{criança} a entrar no mundo dos inventos para compreender fenômenos mais complexos e na própria saga da tradição científica, trabalhando com objetos, plantas e animais.
No que concerne a relação com os números, é fundamental aproveitar a familiaridade da \textbf{criança} com as quantidades e as dimensões, assim como as suas habilidades para tirar e adicionar.
O uso de materiais e de experiências cotidianas com o cálculo para que a constatação e a abstração se encontrem e facilitem o acesso às competências matemáticas iniciais e a representação dos seus símbolos.
Seus movimentos no espaço permitem, também, experimentar e experienciar de forma reflexiva conceitos geométricos.
Inserções de jogos de mesa e jogos eletrônicos abrem possibilidades de desenvolvimento lógico num mundo contemporâneo eivado de constantes desafios neste campo dos inventos que implicam jogos e \textbf{brincadeiras} para \textbf{crianças}.
Nesse particular, é fundamental cuidar das adequações cognitivas e éticas desses jogos.
Assim, o mundo dos números e dos cálculos pode ser experimentado e compreensões podem ser construídas de forma \textbf{lúdica} e processual, incluindo nessa itinerância compreensiva os fenômenos presentes na Natureza, na biologia humana e de outros seres.

% matched lemmas: acolher, adulto, brincadeira, brincar, colega, criança, cuidado, curiosidade, equilíbrio, escuta, gesto, infantil, interessar, lúdico, mímico, pequeno, progressivamente, prédio, som
\end{textsample}
